\slajdid{04-s020}
\begin{frame}[label=04-s020]
\gusttekst
\setlength{\abovedisplayskip}{5pt}\setlength{\belowdisplayskip}{5pt}
Ako postoji razlomljeni deo broja

Sve isto osim što se grupisanje radi:\\
\hspace*{5mm}za celobrojnu vrednost na levo od tačke\\
\hspace*{5mm}za razlomljenu vrednost na desno od tačke

PRIMER\qquad $1A.3C_{16}={?}_2$

Svaku cifru za sebe\qquad $1A.3C_{16}=0001\ 1010\ .\ 0011\ 1100_2$

Pa možemo da izbrišemo vodeće i prateće nule
\[=\textcolor{DEcrvena}{000}1\ 1010\ .\ 0011\ 11\textcolor{DEcrvena}{00}_2=11010.001111_2\]
PRIMER\qquad $101100.1011011_2={?}_{16}$

Da bi olakšali posao pravimo grupe od po četiri cifre
\[10\ 1100\ .\ 1011\ 011_2=\textcolor{DEcrvena}{00}10\ 1100\ .\ 1011\ 011\textcolor{DEcrvena}{0}_2\]
Dodali smo vodeće i prateće nule da bi dobili grupe od četiri cifre
\[\textcolor{DEcrvena}{00}10\ 1100\ .\ 1011\ 011\textcolor{DEcrvena}{0}_2=2C.B6_{16}\]
Potpuno identično binarni u oktalni i obrnuto samo su nam grupe sa po 3 cifre
\end{frame}
